% ======================================================================
%  CAPÍTULO 18 — SOLUCIONES (parte 2: capítulos 1-8)
% ======================================================================

\subsection*{Capítulo 1 --- La amenaza cuántica y el giro
post-cuántico (capítulo~\ref{ch:pqc})}

\begin{enumerate}[leftmargin=1.9em,itemsep=4pt]
  \item Grover busca la clave en $2^{k/2}$ evaluaciones: AES-256
        cae a $2^{128}$ (dentro del listón) y AES-128 a $2^{64}$
        (muy por debajo). La aceleración cuántica es cuadrática en el
        exponente, no total: basta doblar el tamaño de clave.
  \item \emph{Criptogramas}, de menor a mayor: \texttt{mceliece460896}
        ($156$~B) $<$ ML-KEM-768 ($1088$~B) $<$ HQC-128
        ($\approx4497$~B). \emph{Claves públicas}: ML-KEM-768
        ($1184$~B) $<$ HQC-128 ($\approx2249$~B) $<$
        \texttt{mceliece460896} ($\num{524160}$~B). McEliece
        intercambia los extremos: el criptograma más pequeño y la
        clave más grande.
  \item Porque el tráfico cifrado \emph{hoy} ya se graba: cuando
        exista el ordenador cuántico, se descifrará retroactivamente.
        Los datos de vida útil larga están expuestos desde el momento
        en que se capturan; migrar después no repara lo ya grabado, de
        ahí actuar ahora (híbridos incluidos).
  \item \textnormal{(}Investigación.\textnormal{)} La nota del BSI
        (1 de octubre de 2026, \cite{bsi2026}) sostiene, en esencia:
        (a)~los métodos de agosto--octubre (distinguidores de
        $2^{120}$--$2^{124}$; recuperaciones estimadas de
        $2^{140}$--$2^{148}$; el cálculo de una sola ejecución de
        Weis, $2^{99}$--$2^{102}$) estrechan el margen lo bastante
        para tomarse en serio, aunque nada esté roto; (b)~principio de
        precaución: no iniciar despliegues nuevos, preferir ML-KEM,
        FrodoKEM y HQC, y conservar la mitad clásica de los híbridos
        existentes, revisando la TR-02102-1. (Respuesta abierta: lo
        evaluable es distinguir lo demostrado de lo estimado y leer
        la recomendación como decisión de política.)
\end{enumerate}

\subsection*{Capítulo 2 --- Códigos lineales
(capítulo~\ref{ch:lineales})}

\begin{enumerate}[leftmargin=1.9em,itemsep=4pt]
  \item Si $\vx$ y $\vz$ difieren en la posición $j$, entonces
        $\vx_j\neq\mathbf{y}_j$ o $\mathbf{y}_j\neq\vz_j$ (si ambas igualdades se
        dieran, $\vx_j=\vz_j$). Así
        $\supp(\vx+\vz)\subseteq\supp(\vx+\mathbf{y})\cup\supp(\mathbf{y}+\vz)$,
        y tomando cardinales se obtiene la desigualdad triangular.
  \item Una base: las filas $1100$ y $0011$ sirven como $G$ y también
        como $\Hmat$ (las comprobaciones son $x_1+x_2=0$ y
        $x_3+x_4=0$). Las palabras no nulas son $1100$, $0011$ y
        $1111$, de pesos $2$, $2$ y $4$: $d=2$.
  \item Para $y=0101010$: síndrome $=\Hmat_2+\Hmat_4+\Hmat_6
        =011+101+010=100=\Hmat_5$: se corrige la posición $5$
        ($0$-indexada $4$): $\hat y=0101110$, mensaje $0101$. Para
        $y'=1100110$: síndrome $=\Hmat_1+\Hmat_2+\Hmat_5+\Hmat_6
        =110+011+100+010=011=\Hmat_2$: se corrige la posición $2$:
        $\hat y'=1000110$, mensaje $1000$.
  \item Hay $\binom{n}{t}$ palabras de peso $t$: $\binom71=7$ para
        el Hamming; $\binom{3488}{64}\approx2^{456{,}3}$ para
        \texttt{348864}: unos $10^{137}$ candidatos.
  \item Sea $e''$ de peso mínimo con $\Syn(e'')=\Syn(y)$
        (existe: al menos $e$). Entonces $y+e''$ es un codeword (su
        síndrome es $0$) y todo codeword $c$ cumple
        $\wt(y+c)\ge\wt(e'')$, porque $y+c$ tiene el mismo síndrome
        que $e''$ y $e''$ es el mínimo. Así el codeword más cercano
        dista de $y$ exactamente $\wt(e'')$, que es a lo sumo
        $\wt(e)$.
  \item Un código sirve para «publicar $G'=SGP$» si tiene (a)
        decodificación eficiente hasta $t$ (sin la descripción
        secreta no; con ella, Patterson) y (b) ocultación: el disfraz
        hace que $G'$ parezca una matriz aleatoria. Los dos recursos
        de quien descifra son el computacional (Patterson) y el
        algebraico (la descripción $g$ y el orden de $\alpha$, que
        convierten $G'$ en un código con decodificador).
\end{enumerate}

\subsection*{Capítulo 3 --- Cuerpos finitos
(capítulo~\ref{ch:cuerpos})}

\begin{enumerate}[leftmargin=1.9em,itemsep=4pt]
  \item $z^5=z\cdot z^4=z(z+1)=z^2+z$; y
        $z^{-5}=z^{15-5}=z^{10}=z^2+z+1=\hx7$.
  \item $\Tr(1)=1+1+1+1=0$; $\Tr(z)=z+z^2+z^4+z^8
        =z+z^2+(z+1)+(z^2+1)=0$. Tienen traza $1$ exactamente la
        mitad: $8$ elementos (el núcleo de $\Tr$ es un subespacio
        de dimensión $m-1=3$).
  \item $\sqrt{z^3+z^2}$: como $z^3+z^2=z^6$ (pues
        $z^6=z^4z^2=(z+1)z^2=z^3+z^2$), la raíz es $z^3$.
        $\sqrt{z^2+z+1}$: como $z^2+z+1=z^{10}$, la raíz es
        $z^5=z^2+z=\hx6$. En general, $\sqrt{\beta}=\beta^{2^{m-1}}$
        (para $m=4$, elevar al exponente $8$).
  \item $\Tr(\beta^2)=\beta^2+\beta^4+\beta^8+\beta^{16}
        =\beta^2+\beta^4+\beta^8+\beta=\Tr(\beta)$ (en $\FF4$,
        $\beta^{16}=\beta$). La linealidad sobre $\Ftwo$: Frobenius da
        $(\beta+\gamma)^{2^i}=\beta^{2^i}+\gamma^{2^i}$ para cada
        $i$, y sumar conserva la aditividad; $\Tr(a\beta)=a\Tr(\beta)$
        para $a\in\Ftwo$ es inmediato.
  \item $x^2+x+1$ no tiene raíces en $\Ftwo$ ($P(0)=P(1)=1$), luego
        es irreducible; $x^2+x=x(x+1)$ es reducible. El criterio de la
        traza ($x^2+x+a$ irreducible $\iff\Tr(a)=1$) coincide:
        $\Tr(1)=1$ y $\Tr(0)=0$. (Sobre $\FF4$ ambos se parten:
        $x^2+x+1$ tiene las raíces de orden $3$.)
  \item Suma: un XOR por bit. Producto: red de AND+XOR de ancho
        $\approx m^2$. Raíz cuadrada: casi gratis, es $\Ftwo$-lineal
        (una matriz de XORs). Inversa: la más cara (cadena de
        elevar-al-cuadrado y multiplicar, $x\mapsto x^{q-2}$; o
        Itoh--Tsujii; o tablas), tanto por profundidad como por
        dependencia de datos. Patterson la usa $O(t)$ veces.
\end{enumerate}

\subsection*{Capítulo 4 --- Códigos de Goppa
(capítulo~\ref{ch:goppa})}

\begin{enumerate}[leftmargin=1.9em,itemsep=4pt]
  \item Si $g(\alpha_i)=0$, entonces $(x+\alpha_i)$ divide a $g$ y
        $1/(x+\alpha_i)$ no existe módulo $g$ (es un divisor de
        cero); la columna de síndrome
        $\sigma_i=(x+\alpha_i+1)/g(\alpha_i)$ no se puede calcular.
        Equivalente: todas las condiciones se anulan en la posición
        $i$ (el término $c_i/(x+\alpha_i)$ queda libre): aparecerían
        palabras de peso $1$ y la distancia se derrumba.
  \item $A(x)/(x-\alpha_i)=\prod_{j\neq i}(x-\alpha_j)$: evaluado
        en $\alpha_i$ da $\prod_{j\neq i}(\alpha_i-\alpha_j)$. Y
        $A'(x)$ es la suma de los productos que omiten un factor;
        al evaluar en $\alpha_i$, los términos que conservan el
        factor $(x-\alpha_i)$ se anulan, y queda exactamente el
        mismo producto. Es la identidad que traduce la suma de
        fracciones a la fórmula polinomial de la especificación.
  \item Con el ordenamiento natural del juguete ($\alpha_0=0$,
        $\alpha_1=1$): $g(0)=g(1)=z^3$, luego
        $1/g(0)=z^{-3}=z^{12}=\hx f$; la columna $0$ es
        $(\hx f,0)$ y la $1$ es $(\hx f,\hx f)$. (La
        tabla~\ref{tab:lab-columnas} del laboratorio muestra el
        resto.)
  \item $d\ge2t+1$ dice exactamente que la bola de radio $2t$
        alrededor de $0$ solo contiene al codeword nulo: ningún peso
        $\le2t$. No contradice $d>2t+1$ porque es una cota
        \emph{inferior}: el lema garantiza al menos $2t+1$; la
        distancia real puede ser mayor (en el juguete es $5=2t+1$,
        pero la cota no prohíbe $6$ o más en otros códigos).
  \item Con $g=h^2$: la hipótesis de
        \texttt{thm:gopla-2t} (libre de cuadrados) falla, pero la
        respuesta de fondo es la identidad de van Lint
        $\Gamma_2(L,g)=\Gamma_2(L,g^2)$: tomando $g=h$ queda
        $\Gamma_2(L,h^2)=\Gamma_2(L,h)$. El código con $g=h^2$ es
        \emph{el mismo} que con $h$: el «$t$ efectivo» sigue siendo
        $\deg h$, no se gana distancia ni dimensión. McEliece nunca
        lo intentó porque duplicar el grado nominal no cambia el
        código: promesa vacía.
  \item Semillas $\to$ $g$ (con $\textsf{Irreducible}$) y orden de
        $\alpha$ (con $\textsf{FieldOrdering}$) $\to$ el código
        $\Gamma_2(L,g)=\{c:\sum_j c_j/(x+\alpha_j)\equiv0\pmod g\}$
        $\to$ la tabla de síndromes $\tilde{\Hmat}$ ($2\times n$
        sobre $\FF4$, columnas $\alpha_i^j/g(\alpha_i)$) $\to$ la
        matriz binaria $\hat{\Hmat}$ ($8\times16$) $\to$
        eliminación gaussiana $\to$ $(\mathrm{I}_8\,|\,T)$. Se
        publica solo $T$ (junto con $t$); se mantienen secretos $g$,
        el orden $\alpha$ (y las semillas, y $s$ en el KEM). El
        código $\ker(\mathrm{I}_{mt}|T)$ es público.
\end{enumerate}

\subsection*{Capítulo 5 --- Laboratorio 1
(capítulo~\ref{ch:labgoppa})}

\begin{enumerate}[leftmargin=1.9em,itemsep=4pt]
  \item $g(z)=z^2+z+z^3=z^3+z^2+z=\hx e$; $1/\hx e=z^{-11}=z^4
        =z+1=\hx3$; $z\cdot\hx3=z^2+z=\hx6$. La columna $j=2$ de
        la tabla~\ref{tab:lab-columnas} es exactamente $(\hx3,\hx6)$:
        coinciden.
  \item $\tilde{\Hmat}e_5=\tilde{\Hmat}_5=(\hx8,\hx e)$ (pues
        $g(z^2+1)=\hx f$, $1/\hx f=z^3=\hx8$ y
        $(z^2+1)z^3=\hx e$); $\tilde{\Hmat}e_9=(\hx a,\hx5)$ (pues
        $g(z^3+1)=z^6$, $1/z^6=z^9=\hx a$ y
        $(z^3+1)z^9=\hx5$). La suma es
        $(\hx8+\hx a,\hx e+\hx5)=(\hx2,\hx b)$. $e_5+e_9$ se decodifica
        de forma única porque dos palabras de peso $\le2$ con el
        mismo síndrome diferirían en un codeword de peso
        $\le4<d=5$: imposible.
  \item Multiplicando y usando $z^4=z+1$ en cada coeficiente:
        constante $4\cdot\hx a=z^2(z^3+z)=z^3+z^2+z=\hx e$; lineal
        $4\cdot7+9\cdot\hx a=\hx f+\hx5=\hx a$; cuadrático
        $9\cdot7=z^3+z=\hx a$. En el lado derecho,
        $1+(z+z^3)g(x)$ tiene esa misma forma: $z+z^3=\hx a$,
        $1+\hx a z^3=\hx e$. Es el certificado de Patterson.
  \item Con $E=\{j\}$: $\sigma(x)=\prod_{i\in E}(x+\alpha_i)
        =x+\alpha_j$ (grado $1\le t$), y $\sigma'=1$: la ecuación
        clave se satisface con $\sigma=x+\alpha_j$. Patterson
        devuelve el localizador; su raíz $\alpha_j$ señala la
        posición del error.
  \item $(\vx P^{-1})_1=\vx_{\sigma^{-1}(1)}=\vx_{15}=1$ y
        $(\vx P^{-1})_3=\vx_{\sigma^{-1}(3)}=\vx_{13}=0$: la misma
        permutación envía los dos errores de $\vz$ a las posiciones
        $0$ y $10$ ($\supp(\vz P^{-1})=\{\sigma(3),\sigma(12)\}
        =\{10,0\}$). Es exactamente lo que el laboratorio decodifica
        con Patterson.
  \item Dos ataques: (a)~enumerar las $137$ palabras de peso $\le2$
        y quedarse con la que comparte síndrome con $C$
        ($120$ de peso $2$, $16$ de peso $1$, la nula);
        (b)~tabla completa síndrome$\to$palabra mínima ($256$
        entradas, $137$ no triviales): un oráculo de descifrado
        instantáneo. Con $n\approx8000$ el espacio de síndromes pasa
        a $\ge2^{64}$ y $\binom{n}{t}$ es astronómico: mueren la
        tabla y la enumeración, y quedan los algoritmos de la
        parte~\ref{part:seguridad} (ISD, métodos estructurales).
\end{enumerate}

\subsection*{Capítulo 6 --- El McEliece de 1978
(capítulo~\ref{ch:1978})}

\begin{enumerate}[leftmargin=1.9em,itemsep=4pt]
  \item La cadena es
        $\vx P^{-1}=(\vu S)G+\vz P^{-1}$: $\wt(\vz P^{-1})=\wt(\vz)$
        (permutar no cambia el peso), y $\wt(\vz)\le t$ es lo que
        permite aplicar el corolario de corrección al codeword
        $(\vu S)G$. El paso fino: para otro codeword $c'$,
        $\dist(\vx P^{-1},c')\ge d-\wt(\vz)\ge 2t+1-t=t+1>t$:
        el codeword verdadero está \emph{estrictamente} más cerca.
        Con errores de peso $t+1$ «a veces», esa desigualdad se
        rompe (podría haber dos codewords a distancia $t+1$):
        Patterson podría devolver otro codeword y el descifrado
        fallaría o mentiría.
  \item $G'$ para $(1024,524)$: $524\cdot1024=\num{536576}$ bits
        ($\approx65{,}5$~KiB). Para \texttt{348864}:
        $2720\cdot3488=\num{9487360}$ bits
        ($\approx1{,}13$~MB). La parte sistemática
        $(\mathrm{I}_{2720}|T)$ ocupa $mt\lceil k/8\rceil
        =768\cdot340=\num{261120}$~B (la clave pública): el ahorro
        frente a $G'$ completa es $1185920/261120=4{,}54$ veces.
  \item Si $(\vu',\vz')\ne(\vu,\vz)$ diera el mismo $\vx$, entonces
        $(\vu+\vu')G'=\vz+\vz'$: un codeword de peso
        $\le2t<2t+1\le d$, imposible salvo que sea $0$; y si es
        $0$, $\vu=\vu'$ y $\vz=\vz'$.
  \item $\binom{974}{524}/\binom{1024}{524}=2^{-53{,}61}$: es la
        probabilidad de que un conjunto de $k$ coordenadas elegido al
        azar evite las $t=50$ posiciones del error; con
        $2^{53{,}61}$ intentos de coste $k^3\approx2^{27{,}1}$ sale
        $2^{80{,}7}\approx2\cdot10^{24}$ (el «ataque de las $k$
        coordenadas», cotizado en $10^{19}$ en el artículo: el
        recálculo favorece al defensor). $\sum_{i\le50}\binom{1024}{i}
        =2^{284{,}11}$, así que un síndrome al azar se decodifica con
        probabilidad $2^{284{,}11-500}=2^{-215{,}89}$ (el artículo
        decía $2^{-215{,}59}$: coincide en lo esencial).
  \item Porque la NP-dureza es un enunciado \emph{del peor caso},
        \emph{asintótico} y sobre \emph{instancias arbitrarias}:
        (a)~no dice nada de las instancias aleatorias de una familia
        estructurada (los códigos de Goppa); (b)~para $n=3488$ fijo
        no acota nada concreto; (c)~y la seguridad del esquema es una
        conjetura de dureza \emph{en media} sobre su distribución
        (capítulo~\ref{ch:seguridad}).
  \item Con el oráculo que avisa: se envía $\vx\oplus\ve_j$ para cada
        $j$. Si $j$ está en el soporte del error, el peso baja a
        $t-1$ y el descifrado «funciona»; si no, sube a $t+1$ y
        «falla». Bastan $n=1024$ consultas para aprender el soporte
        de $\vz$ y, con él, $\vu$. En el KEM real cada consulta
        devuelve una clave pseudoaleatoria (rechazo implícito): no
        hay aviso que aprender (capítulo~\ref{ch:encapdecap}).
\end{enumerate}

\subsection*{Capítulo 7 --- La reformulación de Niederreiter
(capítulo~\ref{ch:niederreiter})}

\begin{enumerate}[leftmargin=1.9em,itemsep=4pt]
  \item Con $H_0=(\mathrm{I}_{n-k}\,|\,T_0)$, la palabra
        $v=(\gamma,\mathbf{0})$ tiene
        $vH_0^{\top}=\mathrm{I}\cdot\gamma+0=\gamma$. El paso (2) del
        descifrado pide el único $\vf$ de peso $\le t$ con
        $\vf H_0^{\top}=\gamma$: exactamente el problema de
        decodificación del capítulo~\ref{ch:goppa} (encontrar el
        vector de menor peso con un síndrome dado), y Patterson lo
        resuelve operando sobre el polinomio síndrome del vector
        rellenado.
  \item Si $\vf H_0^{\top}=\vf'H_0^{\top}$, entonces
        $(\vf+\vf')H_0^{\top}=0$: $\vf+\vf'$ es un codeword de
        $\Gamma$ de peso $\le2t<d$, luego $\vf+\vf'=0$ y
        $\vf=\vf'$.
  \item La demostración: primero $G'(H')^{\top}
        =S_1GP\,P^{\top}H_0^{\top}S_2^{\top}$: aquí se usa
        $PP^{\top}=\mathrm{I}_n$; después
        $S_1(GH_0^{\top})S_2^{\top}=0$: aquí, la ortogonalidad
        $GH_0^{\top}=0$ (los codewords tienen síndrome nulo).
        Con permutaciones distintas quedaría $P_1P_2^{\top}\ne
        \mathrm{I}_n$ en medio y la ortogonalidad no se seguiría:
        sin ella no hay reducción (el síndrome de $\vx$ ya no es el
        de $\vz$).
  \item Factores $k/(n-k)$: juguete $8/8=1$;
        \texttt{348864}, $2720/768=3{,}54$;
        \texttt{460896}, $3360/1248=2{,}69$;
        \texttt{6688128}, $5024/1664=3{,}02$;
        \texttt{6960119}, $5413/1547=3{,}50$;
        \texttt{8192128}, $6528/1664=3{,}92$.
        El original $(1024,524)$: $k/(n-k)=524/500=1{,}05$: con tasa
        $\approx1/2$, la matriz dual era casi tan grande como la
        generadora; el ahorro, mísero (del orden del $5\%$).
  \item Tamaño del espacio de mensajes: el descifrado debe
        funcionar sobre la imagen de
        $\textsf{FixedWeight}$; admitir pesos menores agrandaría la
        imagen con preimágenes «fáciles» (pesos pequeños, búsqueda
        más ligera) sin ganar entropía. Uniformidad: los teoremas de
        las transformaciones a KEM (capítulo~\ref{ch:kem}) y la
        proposición~\ref{prop:fixedweight-uniform} exigen $\ve$
        uniforme en \emph{el} conjunto de vectores de peso $t$
        exacto; mezclar pesos rompería el modelo y la simetría del
        reto.
  \item Los dos problemas de inversión:
        $\vx=\vu G'+\vz\mapsto\vu$ (McEliece) y
        $\vc=\ve(H')^{\top}\mapsto\ve$ (Niederreiter). Son «el mismo
        problema» por la proposición~\ref{prop:nied-reduccion}: el
        síndrome público de $\vx$ es el de $\vz$ (ortogonalidad
        $G'(H')^{\top}=0$), y recuperar $\ve$ equivale a recuperar
        $\vz$; en ambos casos el fondo es el problema
        de decodificación por síndromes: dada la imagen de un error
        de peso $\le t$, encontrarlo.
\end{enumerate}

\subsection*{Capítulo 8 --- Del cifrado a un KEM
(capítulo~\ref{ch:kem})}

\begin{enumerate}[leftmargin=1.9em,itemsep=4pt]
  \item En la transformación, $\textsf{Encap}$ genera $e$, publica
        $C=\textsf{Enc}(e)$ (aquí $C=(\mathrm{I}_{mt}|T)e$) y deriva
        $K=H(1,e,C)$; $\textsf{Decap}$ calcula $e'=\textsf{Decode}(C)$
        y $K'=H(1,e',C)$. Las dos partes coinciden si y solo si
        $\textsf{Decode}(\textsf{Enc}(e))=e$ con probabilidad $1$: el
        único punto donde entra la corrección de la PKE (el
        $(\varepsilon_{\mathrm{corr}})$ de los teoremas).
  \item Con $K=H(e)$, la consulta al oráculo con la segunda
        codificación $\tilde C$ de $C^*$ (distinta como bytes,
        idéntica como elemento, por el relleno del
        capítulo~\ref{ch:parametros}) devuelve $H(e^*)=K^*$: el reto
        cae con una consulta que no es $C^*$. Con $K=H(1,e,C)$, el
        criptograma entra en el hash: la respuesta del oráculo es
        $H(1,e^*,\tilde C)$, una clave distinta (bindada a la
        codificación exacta) e inútil para el ataque; la etiqueta
        $1$ además separa el dominio de la rama de rechazo
        $H(0,s,C)$.
  \item Devolver $\bot$ obliga a los protocolos a reaccionar
        (alertas, reintentos) y esa reacción es observable: es un
        oráculo de validez. Devolver siempre una clave
        pseudoaleatoria elimina la señal: el atacante no distingue
        «no decodificable» de «decodificable con valor improbable».
        La clave $s$ aparece como secreto de la rama de rechazo: en
        el teorema~\ref{thm:kem-rom} en el término $2q_H/|S|$ (hay
        que adivinar $s$ para simular) y en el
        teorema~\ref{thm:kem-qrom} en los términos de PRF sobre $s$.
  \item Segundo término: $q_D/|P|=2^{64}/2^{456}=2^{-392}$. Tercero:
        $2q_H/|S|=2\cdot2^{64}/2^{256}=2^{-191}$. Para que el tercero
        dejara de ser despreciable haría falta $q_H\approx2^{255}$:
        un número de consultas imposible (y entonces fallarían
        antes los demás términos del modelo).
  \item (a)~Con ventaja $<1/1000$: $2\sqrt{\varepsilon_{\mathrm{OW}}}
        <1/1000$ exige
        $\varepsilon_{\mathrm{OW}}<(1/2000)^2\approx2^{-21{,}9}$.
        (b)~Con ventaja $<2^{-128}$:
        $\varepsilon_{\mathrm{OW}}<2^{-258}$. Es decir: el teorema
        convierte el objetivo CCA de $128$ bits en un requisito
        OW-CPA de $\approx2^{-258}$: los parámetros deben tener
        margen doble, no nulo (es la raíz cuadrada del QROM, que el
        modelo de costes realista suaviza pero no elimina).
  \item La tabla~\ref{tab:kem-cm} lo dice: (1)~sin reencriptar: la
        verificación $\textsf{Enc}(\mathrm{pk},e')=C$ no hace
        falta porque $\textsf{Decode}$ ya exige $\wt(e')=t$ y
        $(\mathrm{I}_{mt}|T)e'=C$; (2)~sin $\mathrm{pk}'$ dentro de
        la clave secreta (no hay reencriptación que lo use);
        (3)~$s$ de $n$ bits en vez de $256$; (4)~$\ve$ de
        $\textsf{FixedWeight}$ en vez de uniforme-genérico. La
        verificación de $\textsf{Decode}$ basta porque, para un
        criptograma, «ser imagen de un $e'$ de peso $t$» y «cumplir
        $\wt(e')=t\ \wedge\ (\mathrm{I}_{mt}|T)e'=C$» son la misma
        condición: la primera se comprueba sin reencriptar.
\end{enumerate}
